\section{Fundamentos de Directional Relational Manifolds}
\label{sec:fundamentos}

Um DRM é uma estrutura geométrica na qual cada ponto $p$ possui um conjunto de
direções ativas $D(p)$. A dimensionalidade local é definida por
\begin{equation}
  \dim_D(p) = |D(p)|,
\end{equation}
ou, em formulações lineares, pelo posto efetivo do conjunto direcional. Ao
contrário de uma variedade Riemanniana clássica, onde a dimensão local é
constante, um DRM admite que $\dim_D(p)$ varie ao longo do espaço. A estrutura
básica é um par $(M,D)$, com $M$ topológico e $D$ uma atribuição mensurável de
direções ativas que varia continuamente em sentido apropriado.

A métrica relacional é definida apenas sobre as direções ativas:
\begin{equation}
  g_p: D(p) \times D(p) \rightarrow \mathbb{R}.
\end{equation}
Quando $D(p)$ é constante, finito e ortonormal, a estrutura reduz ao caso
Riemanniano usual. Quando $D(p)$ tem dimensão constante e $g_p$ coincide com a
métrica de Fisher em uma família estatística, o DRM reduz à geometria da
informação. Dessa forma, DRM pode ser visto como uma generalização que permite
postos variáveis e modos não ortogonais \cite{amari_information_geometry,
muniz_drm_v11}.

No contexto computacional deste trabalho, o conjunto direcional não é
representado explicitamente como uma base simbólica, mas por coordenadas
aprendidas em um manifold de dimensão $d_m$ e por um tensor métrico low-rank.
A parametrização usada no modelo é
\begin{equation}
  G(x) = I + U(x) U(x)^\top,
\end{equation}
onde $U(x) \in \mathbb{R}^{d_m \times r}$ é produzido por uma rede
\emph{MetricNet}. Essa forma garante positividade semidefinida e mantém o custo
de atenção viável, pois a distância métrica pode ser computada por
\begin{equation}
  d_G^2(q,k) = \|\Delta\|^2 + \|U(q)^\top \Delta\|^2,
  \quad \Delta = q-k.
\end{equation}

O teorema de convergência toroidal do DRM original afirma que, sob hipóteses de
limitação, recorrência e estabilidade estrutural, a dinâmica pode apresentar
conjuntos limite compactos com topologia toroidal \cite{muniz_drm_v11}. No
modelo proposto, essa afirmação não é assumida como garantia universal; ela é
tratada como alvo empírico verificável por homologia persistente.
